% Reconstructed after environment loss from the complete accepted conversation proof.
\statement{C.43}{Proved}{Self-centralizing sections for the regular pants involution}
Let $F=\widehat F(x,y)$ and $z=(xy)^{-1}$. Suppose an involution
$\sigma$ interchanges the peripheral classes of $x$ and $y$ with
exponent $-1$, and preserves the class of $z$ with exponent $-1$.
Then $E=F\rtimes_\sigma C_2$ has exactly one $F$-conjugacy class
of sections, and every section $S$ satisfies $C_F(S)=1$.
Every open $\sigma$-stable subgroup inherits self-centralizing
torsion in its split extension. In particular, if $U$ is also normal
in $F$, the map $C_{F/\Phi_2(U)}(\sigma)\to F/U$ has trivial image.

\begin{proof}
Put $L_x=\overline{\langle x\rangle}$, and define $L_y,L_z$
similarly. C.5 gives their geometric normalizers. No odd element of
$E$ normalizes $L_x$, because its outer action interchanges two
distinct peripheral classes. Thus $D_x=N_E(L_x)=L_x$. Its two
$F$-orbits give, with the actual conjugating transport for the second,
\[
 \operatorname{Res}_F(E/D_x)\cong F/L_x\sqcup F/L_y.
\]
On the other hand $D_z=N_E(L_z)=L_z\rtimes_{-1}C_2$: an odd
normalizer lift inverts $L_z$, and its square belongs to $L_z$
and is fixed by inversion, hence is one.

Use one summand per peripheral orbit, with $k=\mathbb F_2$:
\[
 \Delta=\ker\bigl(k[[E/D_x]]\oplus k[[E/D_z]]\longrightarrow k\bigr).
                                                                    \tag{C.43.1}
\]
Its restriction to $F$ is exactly the original three-boundary
surface-pair augmentation module. Therefore the complete proof of
C.6 applies to this orbit-indexed family. The dual fibre rows,
transposed differentials, fundamental-class trace and filtered
perfect pairing are unchanged. For each finite self-dual $E$-module,
localized boundary restriction consequently has half the boundary
dimension. The swapped cusp pair remains in (C.43.1) throughout this
duality argument. Only afterwards does its localized boundary term
vanish, because $D_x\cong\widehat{\mathbb Z}$ has
$2$-cohomological dimension one.

The action on $H_1(F,k)$ is the regular two-dimensional
$k[C_2]$-module. The split Hochschild--Serre sequence gives
$\dim_kH^n(E,k)=1$ for $n\ge2$. The full restriction theorem
\cite[Theorem 2]{CEKochloukovaZalesskii2007}, applied as in
C.37, gives exactly one global section class. This count does not
assume that its centralizer is trivial.

Fix that class $S$ and put $B=C_F(S)$. Its normalizer is
$B\times S$. The dihedral boundary $D_z$ has two local section
classes, represented by $r$ and $zr$, each with normalizer $C_2$;
both map to the single global class. For an arbitrary finite
self-dual permutation $E$-module $M$, high-degree restriction
identifies the absolute term with the cohomology of $B\times S$
and the remaining boundary term with two copies of the cohomology
of $S$. After the actual conjugating coefficient identifications,
the two endpoint maps have the same kernel: each is the restriction
\[
 H^\bullet(B\times S,M)[u^{-1}]_0
       \longrightarrow H^\bullet(S,M)[u^{-1}]_0.                 \tag{C.43.2}
\]
If its target has dimension $d$, the joint endpoint map has the
single restriction rank, at most $d$. The orbit version of C.6
makes that rank half of $2d$, hence $d$. Thus (C.43.2) is onto
for every such $M$.

If $1\ne h\in B$, separate $h$ from the finite subgroup $S$
in a finite quotient $D$ of $E$, and take
$M=k[D/\overline S]$. Positive $S$-cohomology has the fixed cosets
as a basis. The element $h$ moves the identity-coset vector to a
distinct fixed-coset vector. Restriction from $B\times S$, however,
has $B$-invariant image by naturality under inner automorphisms.
This contradicts surjectivity and proves $B=1$.

Every involution of $U\rtimes_\sigma C_2$, for open stable $U$,
is an involution of $E$ and has centralizer exactly $C_2$.
The group $U$ is free profinite of finite rank, so C.42 applies
to every such split extension. For open normal stable $U$, it gives
the claimed first relative-dyadic head comparison.
\end{proof}

\paragraph{The common-normalizer consequence.}
For an actual full negative modular involution $b$, suppose the
chosen literal normalizers $d_s=\beta^{-1}b$ and $d_r=J^{-1}b$
are involutory with the prescribed heads. The two involutions
$sd_s$ and $d_r$ have the same $(s,c)$ head in $S_3\times C_2$.
Its preimage has the regular pants action, so they are $F$-conjugate
by C.37 and their conjugator is unique by C.43. For open normal
stable $U$, any two finite conjugators modulo $\Phi_2(U)$ have
the same $F/U$ image by C.42. This is a necessary full-word
constraint. It does not recognize a free product or establish
generation or the third cusp. The finite $r$-normalizer remains a
variable: A.16 can adjust it when the required base transporter is
available. Failure for one arbitrary stored finite frame is
therefore not a full norm obstruction.
